\section{电磁感应 II}

\subsection{自感和互感}

\subsubsection{自感系数}

对于非铁磁质，通过线圈的总磁通量和通过线圈的电流成正比：
$$
\Psi = L I
$$
这里 $L$ 称为\textbf{自感系数}。当 $I$ 在变化时，线圈中产生的电动势：
$$
\varepsilon = -\dv{\Phi}{t} = -L \dv{I}{t}
$$

\subsubsection{互感系数}

对于两个线圈 $L_1, L_2$，记 $L_1$ 产生的磁场通过 $L_2$ 的磁通量为 $\Psi_{21}$。若介质非铁磁质，有
$$
\Psi_{21} = M_{21} I_1
$$
这里 $M_{21}$ 称为\textbf{$L_2$ 对 $L_1$ 的互感系数}。当 $I_1$ 变化时，$I_2$ 中的 $\Psi_{21}$ 也发生变化，产生感应电动势：
$$
\varepsilon_{21} = - \dv{\Psi}{t} = -M_{21} \dv{I_1}{t}
$$
一般情况下，都有 $M_{12} = M_{21}$。

\subsection{磁场能量}

\subsubsection{自感能}

\textbf{自感能}定义为在线圈自感作用下建立电流 $I$ 所做的功：
$$
A = \int_{0}^{I} L i \dd{i} = \frac{1}{2} L I^2
$$

\subsubsection{互感能}

\textbf{互感能}定义为在线圈互感作用下建立电流 $I$ 所做的功。对于一对电流 $I_1, I_2$，互感能为：
$$
A = M_{12} I_1 I_2
$$

一般系统储存的总磁场能为自感能和互感能的和：
$$
W = \frac{1}{2} \sum_{i=1}^{n} L_i I_i^2 + \sum_{i < j} M_{ij} I_i I_j
$$

\subsubsection{磁场能和磁能密度}

\begin{definition}[磁能密度]
	磁场的\textbf{磁能密度}定义为单位体积的磁场能量，即：
	$$
	w_m = \dv{W_m}{V}
	$$
\end{definition}

对于各种磁场均有
$$
w_m = \frac{1}{2} \vect{B} \vdot \vect{H}
$$

对于电磁场，我们有
$$
w = w_e + w_m = \frac{1}{2} \ab(\vect{D} \vdot \vect{E} + \vect{B} \vdot \vect{H})
$$

\subsection{麦克斯韦电磁场理论}

\subsubsection{位移电流}

在非稳恒电流情形下，原来的安培环路定理不再普遍适用。为解决电容器充电过程中回路积分与取面方式矛盾的问题，麦克斯韦引入了\textbf{位移电流}。

\begin{definition}[位移电流]
	通过某个曲面 $S$ 的位移电流定义为电位移通量对时间的变化率：
	$$
	I_d = \dv{t} \iint_S \vect{D} \vdot \dd{\vect{S}}
	$$
\end{definition}

对应的位移电流密度为
$$
\vect{j}_d = \pdv{\vect{D}}{t}
$$

位移电流的实质是\textbf{变化着的电场}，它和传导电流一样能够激发磁场。

\subsubsection{全电流和安培环路定理的普遍形式}

定义\textbf{全电流}为传导电流与位移电流之和：
$$
I_{\text{total}} = I_c + I_d
$$

于是安培环路定理可推广为
$$
\oint_L \vect{H} \vdot \dd{\vect{l}} = \iint_S \ab(\vect{j} + \pdv{\vect{D}}{t}) \vdot \dd{\vect{S}}
$$

即磁场强度沿任一闭合回路的线积分，在数值上等于该回路所围曲面内传导电流和位移电流的代数和。

\subsubsection{麦克斯韦方程组}

\textbf{积分形式}：
$$
\begin{dcases}
	\oiint_S \vect{D} \vdot \dd{\vect{S}} = \iiint_V \rho \dd{V} \\
	\oint_L \vect{E} \vdot \dd{\vect{l}} = - \dv{t} \iint_S \vect{B} \vdot \dd{\vect{S}} \\
	\oiint_S \vect{B} \vdot \dd{\vect{S}} = 0 \\
	\oint_L \vect{H} \vdot \dd{\vect{l}} = \iint_S \ab(\vect{j} + \pdv{\vect{D}}{t}) \vdot \dd{\vect{S}}
\end{dcases}
$$

\textbf{微分形式}：
$$
\begin{dcases}
	\div \vect{D} = \rho \\
	\curl \vect{E} = - \pdv{\vect{B}}{t} \\
	\div \vect{B} = 0 \\
	\curl \vect{H} = \vect{j} + \pdv{\vect{D}}{t}
\end{dcases}
$$

在线性各向同性介质中，还需要配合介质方程
$$
\vect{D} = \varepsilon \vect{E},\quad \vect{B} = \mu \vect{H}
$$
才能确定电磁场。

麦克斯韦方程组的意义在于：

\begin{enumerate}
	\item 麦克斯韦方程组概括了电磁场的基本规律，完善了宏观电磁场理论；
	\item 方程组在任何惯性系中形式相同，具有洛伦兹变换不变性；
	\item 方程组预言了电磁波的存在，在真空中的传播速度为
	$$
	c = \frac{1}{\sqrt{\mu_0 \varepsilon_0}}
	$$
	其数值与光速相同，因此光是电磁波。
\end{enumerate}

\subsection{作业}

\begin{problem}
	习题 6.7
	\begin{solution}
		导体棒下方的回路面积为
		$$
		S_2 = \frac{\pi R^2}{6} - \frac{\sqrt{3}}{4} R^2 = \ab(\frac{\pi}{6} - \frac{\sqrt{3}}{4}) R^2
		$$
		导体棒上方的回路面积为
		$$
		S_1 = \pi R^2 - S_2 = \ab(\frac{5 \pi}{6} + \frac{\sqrt{3}}{4}) R^2
		$$
		那么得到
		$$
		U_{DA} = k S_1 - k S_2 = k \ab(\frac{2\pi}{3} + \frac{\sqrt{3}}{2}) R^2
		$$
	\end{solution}
\end{problem}

\begin{problem}
	习题 6.9
	\begin{solution}
		\begin{enumerate}
			\item[\textbf{1)}] 中心小线圈内部，磁场近似均匀，因此：
			$$
			B = 2\pi R N \cdot \frac{\mu_0}{4 \pi} \frac{I}{R^2} = \frac{\mu_0 N I}{2 R},\ \Phi = N' B S = \frac{\mu_0 N N' I S}{2 R}
			$$
			所以互感系数 $M = \frac{\Phi}{I} = \frac{\mu_0 N N' S}{2 R} \approx 6.28 \times 10^{-6}\,\mathrm{H}$。

			\item[\textbf{2)}] 小线圈中感应电动势：
			$$
			\varepsilon = \dv{\Phi}{t} = M \cdot \dv{I}{t} = 3.14 \times 10^{-4} \,\mathrm{V}
			$$
		\end{enumerate}
	\end{solution}
\end{problem}

\begin{problem}
	习题 6.10
	\begin{solution}
		由题知，互感系数为
		$$
		M = \frac{L - L'}{4} = 0.30 \,\mathrm{H}
		$$
		解方程：
		$$
		\begin{dcases}
			L_1 + L_2 + 2 M = 1.90 \,\mathrm{H} \\
			\frac{M}{\sqrt{L_1 L_2}} = 0.50
		\end{dcases}
		$$
		最后解得 $L_1 = 0.90\,\mathrm{H},\ L_2 = 0.40\,\mathrm{H}$。
	\end{solution}
\end{problem}

\begin{problem}
	习题 6.12
	\begin{proof}
		分析知，导线距离圆心 $r$ 处的磁场都相等，根据安培环路定理：
		$$
		2\pi r \cdot B = \mu_0 I \frac{r^2}{R^2} \implies B = \frac{\mu_0 I r}{2 \pi R^2}
		$$
		因此磁能密度为
		$$
		w_m = \frac{1}{2} B H = \frac{1}{2 \mu_0} B^2 = \frac{\mu_0 I^2 r^2}{8 \pi^2 R^4}
		$$
		因此单位长度的磁场能量为
		$$
		W = \int_{0}^{R} \frac{\mu_0 I^2 r^2}{4 \pi^2 R^4} \cdot 2 \pi r \dd{r} = \frac{\mu_0 I^2}{4 \pi R^{4}} \ab[\frac{1}{4} r^{4}]_{0}^{R} = \frac{\mu_0 I^2}{16 \pi}
		$$
	\end{proof}
\end{problem}

\begin{problem}
	习题 6.13
	\begin{solution}
		\begin{enumerate}
			\item[\textbf{1)}] 使用安培环路定理，在距离轴心 $r$ 处，磁感应强度为：
			$$
			B(r) = \begin{dcases}
				\frac{\mu_0 I r}{2 \pi a^2} &,\ 0 < r < a \\
				\frac{\mu_0 I}{2 \pi r} &,\ a \le r < b \\
				\frac{\mu_0 I \ab(c^2 - r^2)}{2 \pi r \ab(c^2 - b^2)} &,\ b \le r < c \\
				0 &,\ r \ge c
			\end{dcases}
			$$
			因此，根据 $w_m = \frac{1}{2 \mu_0} B^2$ 得到：
			\begin{itemize}
				\item 在导线内：$\displaystyle w_m = \frac{\mu_0 I^2 r^2}{8 \pi^2 a^{4}}$；
				\item 在导线和圆筒之间：$\displaystyle w_m = \frac{\mu_0 I^2}{8 \pi^2 r^2}$；
				\item 在圆筒内：$\displaystyle w_m = \frac{\mu_0 I^2 \ab(c^2 - r^2)^2}{8 \pi^2 r^2 \ab(c^2 - b^2)^2}$；
				\item 在圆筒外：$w_m = 0$。
			\end{itemize}

			\item[\textbf{2)}]
			$$
			\begin{aligned}
				W & = \int_{0}^{a} \frac{\mu_0 I^2 r^2}{8 \pi^2 a^{4}} \cdot 2 \pi r \dd{r} + \int_{a}^{b} \frac{\mu_0 I^2}{8 \pi^2 r^2} \cdot 2 \pi r \dd{r} + \int_{b}^{c} \frac{\mu_0 I^2 \ab(c^2 - r^2)^2}{8 \pi^2 r^2 \ab(c^2 - b^2)^2} \cdot 2 \pi r \dd{r} \\
				& = \frac{\mu_0 I}{16 \pi} + \frac{\mu_0 I^2}{4\pi} \ln \frac{b}{a} + \frac{\mu_0 I^2}{4\pi \ab(c^2 - b^2)^2} \ab(c^{4} \ln \frac{c}{b} - c^2 \ab(c^2 - b^2) + \frac{c^{4} - b^{4}}{4}) \\
				& \approx 1.72 \times 10^{-5} \,\mathrm{J / m}
			\end{aligned}
			$$
		\end{enumerate}
	\end{solution}
\end{problem}